\section{RNN 的训练：时间反向传播}

\subsection{计算图与损失函数}

在 many-to-many 设置中，RNN 在每个时间步都产生输出并计算损失。总损失为各时间步损失之和：

$$L = \sum_{t=1}^{T} L_t$$

\begin{figure}[H]
\centering
\includegraphics[width=0.95\textwidth]{slides-images/slide-017.jpg}
\caption{RNN 的计算图：展示了输入、隐藏状态、输出和损失之间的依赖关系。注意所有时间步共享相同的权重 $W$\protect\footnotemark}
\end{figure}
\footnotetext{Slides 第18页。}

\begin{importantbox}{时间反向传播（BPTT）}
由于所有时间步共享相同的权重矩阵 $W$，计算 $\frac{\partial L}{\partial W}$ 时需要：
\begin{enumerate}
    \item 分别计算每个时间步对 $W$ 的梯度贡献
    \item 将所有时间步的梯度\textbf{累加}
\end{enumerate}
这个过程称为\textbf{Backpropagation Through Time (BPTT)}。可以将其理解为：在概念上将 $W$ 视为每个时间步使用了不同的``副本''，分别计算各副本的梯度，最后因为它们实际上是同一个 $W$ 而将梯度求和。
\end{importantbox}

\subsection{BPTT 的内存问题}

对于长序列，BPTT 需要在内存中保存所有时间步的激活值和梯度。这会导致 GPU 内存迅速耗尽。

\begin{figure}[H]
\centering
\includegraphics[width=0.95\textwidth]{slides-images/slide-019.jpg}
\caption{标准 BPTT 需要将整个序列的计算图保存在内存中\protect\footnotemark}
\end{figure}
\footnotetext{Slides 第20页。}

\subsection{截断 BPTT}

\textbf{解决方案}：将长序列分割为固定长度的``块''（chunk），在每个块内独立执行前向传播和反向传播。

\begin{figure}[H]
\centering
\includegraphics[width=0.95\textwidth]{slides-images/slide-020.jpg}
\caption{截断 BPTT（Truncated BPTT）：将序列分块处理。每个块的初始隐藏状态来自上一块的最终隐藏状态，但梯度不再跨块传播\protect\footnotemark}
\end{figure}
\footnotetext{Slides 第21页。}

具体流程：
\begin{enumerate}
    \item 处理第一个块：从 $h_0$ 开始，计算隐藏状态和损失，执行反向传播，更新权重
    \item 处理第二个块：用第一个块的最终隐藏状态作为起始 $h_0$，但\textbf{不传播梯度}
    \item 重复直到序列结束
\end{enumerate}

\begin{warningbox}{截断 BPTT 的信息损失}
截断 BPTT 是一种近似方法：
\begin{itemize}
    \item \textbf{前向传播}：隐藏状态在块之间传递，信息得以保留
    \item \textbf{反向传播}：梯度在块边界被截断，长距离依赖的梯度信号会丢失
\end{itemize}
在 many-to-many 设置中，信息损失更为明显，因为每个块只看到局部的损失。理想情况下当然是将整个序列放入内存一次性计算，但这对长序列不可行。
\end{warningbox}

\subsection{本章小结}

RNN 的训练通过 BPTT 实现，本质上是在展开的计算图上执行标准反向传播。截断 BPTT 通过将序列分块来控制内存消耗，但代价是丢失了跨块的梯度信息。

%% ========== 字符级语言模型 ==========

